\section{Introduction}\label{sec:introduction}

For $f\in L^1_{\mathrm{loc}}(\R^2)$, the centered disk maximal function is
\[
 Mf(x)=\sup_{r>0}\frac1{\pi r^2}\int_{B(x,r)}|f(y)|\,\dd y,
 \qquad B(x,r)=\{y\in\R^2:|y-x|<r\}.
\]
Every averaging disk is centered at the point of evaluation, and all
positive radii are allowed. Kinnunen proved that this operator is
bounded on $W^{1,p}$ for $1<p<\infty$
\cite[Theorem~1.4]{Kinnunen1997}. The pointwise control of its gradient
by $M(|\nabla f|)$ gives the corresponding norm bound through the
strong $L^p$ maximal inequality. At $p=1$ that last step fails.
Haj{\l}asz and Onninen asked whether the gradient estimate itself
nevertheless survives \cite[Question~1]{HajlaszOnninen2004}.
We give a positive resolution of the planar centered-disk case of
their endpoint question.

\begin{theorem}\label{thm:main}
There is an absolute constant $C<\infty$ such that, for every
real-valued $f\in W^{1,1}(\R^2)$, the function $Mf$ is locally
integrable, belongs to $W^{1,1}_{\mathrm{loc}}(\R^2)$, and satisfies
\begin{equation}\label{eq:main}
 \int_{\R^2}|\nabla Mf|\,\dd x
 \le C\int_{\R^2}|\nabla f|\,\dd x.
\end{equation}
In particular, its distributional derivatives are represented by
globally integrable functions.
\end{theorem}

Here $W^{1,1}(\R^2)$ is the inhomogeneous Sobolev space: the input and
its weak first derivatives are globally integrable. Vector norms are
Euclidean. No radiality or compact-support condition is imposed on
the input. The maximal function itself need not be globally
integrable, even for a nonzero nonnegative smooth function of compact
support. The constant in \eqref{eq:main} is finite and absolute; its
optimal value is not determined. The argument below concerns disks
in the plane.

\subsection{The signed estimate behind the theorem}

The intermediate estimate retains the sign of the input. For real
$g\in C_c^\infty(\R^2)$, put
\[
 A_tg(x)=\frac1{\pi t^2}\int_{B(x,t)}g(y)\,\dd y,
 \qquad S_{a,b}g(x)=\sup_{a\le t\le b}A_tg(x).
\]
For fixed $0<a\le b<\infty$, this envelope is compactly supported and
Lipschitz. The issue is to bound its gradient uniformly as the range
of radii grows.

\begin{proposition}[Signed finite-band estimate]\label{prop:finite-band}
There is an absolute constant $C$ such that, for every real
$g\in C_c^\infty(\R^2)$ and $0<a\le b$ with $b/a$ a power of two, including one,
\begin{equation}\label{eq:intro-band}
 \int_{\R^2}|\nabla S_{a,b}g|\,\dd x
 \le C\int_{\R^2}|\nabla g|\,\dd x.
\end{equation}
The constant is independent of $a$, $b$, and the number of dyadic
radius scales between them.
\end{proposition}

The proof of Proposition~\ref{prop:finite-band} is given in Section~\ref{sec:assembly}, after the local and multiscale estimates described below. A single radius, $a=b$, is covered directly by convolution. The signed form of Proposition~\ref{prop:finite-band} has a specific purpose in the final regularity argument.
After extension to Sobolev inputs, one applies it to
$g_Q=\chi_Q(g-c_Q)$, where $c_Q$ is a local mean and $\chi_Q$ equals
one on the small disks being tested. On those disks,
$A_tg_Q=A_tg-c_Q$, so the corresponding envelopes have the same
gradient. Poincar\'e's inequality then bounds $\|\nabla g_Q\|_1$
by the gradient of $g$ on a slightly larger square. The function
$g_Q$ can take both signs. Signed averages also occur in Weigt's
dyadic and uncentered cube estimates
\cite{Weigt2023Dyadic,Weigt2025Cubes}; their role here is to permit
this local subtraction.

\subsection{Earlier endpoint results}

In one dimension, Tanaka proved an $L^1$ derivative bound for the
uncentered maximal operator \cite[Theorem~1]{Tanaka2002}.
Aldaz and P\'erez L\'azaro extended the result to inputs of bounded
variation, established absolute continuity of the maximal function,
and obtained the sharp variation constant $1$
\cite[Theorem~2.5]{AldazPerezLazaro2007}.
Kurka proved a variation bound with an absolute constant for the
centered operator on functions of bounded variation, together with
local Sobolev regularity and an integrable derivative for Sobolev
inputs \cite[Theorem~1.2 and Corollary~1.4]{Kurka2015}.
The ordering of intervals used in these arguments does not directly
organize a family of centered disks.

In higher dimensions, Aldaz and P\'erez L\'azaro proved regularity for
block-decreasing inputs, invariant under coordinate sign changes and
nonincreasing in each positive coordinate, and uncentered maximal operators over balls of
unconditional norms; their weak-gradient conclusion applies, in particular,
to block-decreasing $W^{1,1}$ input and Euclidean balls
\cite[Theorem~11]{AldazPerezLazaro2009Block}.
Luiro proved the $L^1$ gradient estimate for
radial Sobolev input and the uncentered ball operator
\cite[Theorem~3.12]{Luiro2018Radial}.
Weigt established variation bounds for the dyadic maximal operator
\cite{Weigt2023Dyadic} and for the uncentered maximal operator over
axis-parallel cubes \cite[Theorem~1.1]{Weigt2025Cubes}.
Their geometry and indexing differ from centered disks. Moreover,
a variation estimate controls a derivative measure, which may still
have a singular part. For fractional orders $0<\beta<1$, Beltran and Madrid established
centered endpoint gradient bounds for radial $W^{1,1}$ input
\cite[Theorem~1.2]{BeltranMadrid2020CenteredFractional}.
Weigt subsequently proved endpoint gradient estimates for both
the centered and uncentered operators
\cite[Theorem~1.1]{Weigt2022Fractional}. The positive order restricts
overlap among optimizing balls in that proof; the mechanism
degenerates at order zero.

Lahti and Weigt's conditional regularity theorem upgrades local
bounded variation of the centered maximal function to local Sobolev
regularity when the input is Sobolev
\cite[Theorem~1.2]{LahtiWeigt2025}. It assumes the requisite variation
bound. We prove that bound and then remove the singular part directly
using the signed localization above. Finally, boundedness and
continuity of the gradient map are distinct questions. The
one-dimensional continuity results of Carneiro, Madrid and Pierce
for the uncentered operator \cite{CarneiroMadridPierce2017}, and of
Gonz\'alez-Riquelme for the centered operator
\cite{GonzalezRiquelme2023Continuity}, concern this further issue.
Theorem~\ref{thm:main} asserts the planar norm estimate.

\subsection{How the finite-band estimate is proved}

Write
\begin{equation}\label{eq:signed-average}
 u_t=A_tg,\qquad h=(\nabla g)\,\dd x,\qquad
 V=|h|(\R^2)=\int|\nabla g|.
\end{equation} Scaling reduces
Proposition~\ref{prop:finite-band} to radii in $[2^{-m},1]$.
The proof has one geometric input and two separate multiscale
constructions. The first construction organizes scalar cancellation;
the second chooses directions and pays for their changes in both
scale and space.

\paragraph{Disk moments and concentration in a shell.}
At an interior maximizing radius, the radial derivative of $u_t$
vanishes. A rotational moment vanishes on every disk. In complex
notation these identities give a zero first moment, which is
preserved by suitable positive mixtures of nearby disks. Variations
of optimizing balls also produce weighted-gradient identities in
Luiro's work \cite[Lemmas~2.2 and~2.4]{Luiro2018Radial}; we derive
the identities needed for centered signed averages directly.

Identify $\R^2$ with $\mathbb C$, let $x$ be the disk center, and let
$z$ be a point of its thin boundary shell. The purpose of the moment
identity is visible in the affine factor
\[
 W_z(y)=1-\frac{\overline{y-x}}{\overline{z-x}}.
\]
This factor vanishes at $z$. Inserting it in a difference of balanced
disk kernels therefore makes mass near $z$ inexpensive. This leads
to the geometric question: when is almost all relevant shell mass
concentrated near one point? A sweep by nested disks tangent at that
point answers it outside one exceptional set of centers, uniformly
over every radius in the band. Crucially, it requires a bound on
positive mass only in an inner disk; the error mass is controlled
separately.

These observations yield a local directional estimate. On a square
of side $r$, decompose $h$ into positive measures with directions
in one cone and localized errors. Errors smaller than the
observation scale have zero total vector mass; errors at scale
$r$ are charged their full variation. An
error of variation $b$ supported at scale $s\le r$ is charged
$b(s/r)^\alpha$, with a fixed $\alpha=1/100$. The resulting sum
controls the integral over the square of the supremum, over locally
maximizing radii, of the gradient component pointing against the
cone. Companion estimates handle band endpoints and the value
oscillation
\[
 v_k(x)=\sup_{0<t,t'\le4r_k}|u_t(x)-u_{t'}(x)|,
 \qquad r_k=2^{-k}.
\]
Those estimates include the endpoint arguments and the comparison
of successively narrower kernels; shell concentration alone does
not supply them.

\paragraph{Scalar cancellation and a common price for each cell.}
Let $\phi_q$ be the product of the two one-dimensional tents at a
vertex $q$ of the level-$k$ dyadic grid: each factor equals one at its
vertex coordinate and decreases linearly to zero one mesh width away. Set $h_q=\int\phi_q\,\dd h$, $a_q=|h_q|$.
The tents form positive partitions of unity and refine positively
onto finer grids. We express vectors in each of finitely many bases formed by the two
boundary directions of a cone. For each scalar coordinate, we start
with the actual signed density at the finest level.
Pair its positive and negative parts in equal masses, retain the
constant-sign remainder, and repeat this operation on the weighted
remainders at each parent. This bottom-up construction retains
the support and birth level of every pair, including cancellation
inside the finest tents. Its identity is an identity of measures,
so there is no unrepresented terminal density.

The masses paired at successive levels form a telescoping sum.
Charging each pair at coarser
observation scales with the factor $(s/r)^\alpha$ gives a
nonnegative price $D_Q$ for every dyadic cell, with
\begin{equation}\label{eq:overview-price-budget}
 \sum_Q D_Q\le CV.
\end{equation}
One price works for every direction in the fixed finite palette.
The local estimates then cost $D_Q$ together with appropriate
nearby vector masses $a_q$ at the observation scale. Multiscale oscillation budgets already
play a role in Kurka's one-dimensional proof
\cite[Lemma~1.5]{Kurka2015}; the scalar measure decomposition and
its common prices are proved here.

\paragraph{Vector particles and the cost of choosing directions.}
The remaining vector masses require a different construction.
Vector refinement makes their total increase with the level,
while keeping it at most $V$. We realize these masses by a finite
positive measure on refinement paths, adding mass when needed
and recording its birth level. This construction preserves the
complete path, including levels before birth, so that it can
compare choices of direction made from the first scale onward.
These vector particles differ from the scalar cancelling pairs:
their role is to turn estimates along paths into sums of vector-mass
charges over cells.

Two costs must be controlled. A change of direction starts a new
run of radius scales and creates endpoint terms. Spatial variation
of the choice of direction creates derivatives of its weights
when we integrate by parts. Section~\ref{sec:labels} constructs a finite-state
law at each point $x$ and bounds sums of both costs, weighted
by these vector masses, by $CV$, with the spatial suprema needed
for the cell estimates. Vector-additivity
defects control changes along particles, and comparison at nearby
points transfers that control to the spatially varying law.

\paragraph{Assembly and cancellation of differentiated rows.}
After the common prices and the direction costs have been
bounded, we assemble the radius bands. A run is a maximal consecutive
interval of scales carrying the same direction. For each history,
apply the local directional estimate to the supremum over the
radii of every run, then average over histories. The common
prices and the bounds for direction changes pay for the
unfavorable components and the endpoints of runs.

For each run, subtract the average at its starting scale from
its envelope; the endpoint estimates pay for the gradients of
these baselines. The vector quantity to be integrated by parts
is the sum of the resulting envelopes multiplied by their run
directions. Its history weight depends on $x$. If
$\pi_k(x;n,n')$ is the probability of passing from direction $n$
to direction $n'$ at step $k$, normalization of each row gives
\[
 \sum_{n'}\nabla_x\pi_k(x;n,n')=0.
\]
Before taking absolute values, we can therefore subtract the
analogous vector sum for the preceding history, ending its last
run at level $k-1$. That sum is independent of the next direction.
The difference contains only extensions of the last run and new
runs, and is bounded by $C\sum_{l\ge k}v_l$. The local value
estimate contributes the scale $r_l$, while a differentiated
transition row has scale $1/r_k$. Their ratio $r_l/r_k$ is
exactly the geometric weight controlled by the direction-cost
estimate. Together with the common-price budget, these bounds
give \eqref{eq:intro-band}.

\subsection{From the finite-band estimate to a weak gradient}

Approximation extends Proposition~\ref{prop:finite-band} to signed Sobolev
inputs. Applying it to $g=|f|$ and increasing the radius bands first
gives a finite vector Radon measure representing the derivative
of $Mf$. To prove that this measure has no singular part, fix a
compact null set and a positive radius threshold. The envelope
over larger radii is Lipschitz, so its contribution vanishes
as the neighborhood of the null set shrinks. For small radii, the signed
cutoff and local mean subtraction described above reduce the cost
to the integral of $|\nabla g|$ within a fixed multiple of the
threshold distance. After shrinking the neighborhood, let the
radius threshold tend to zero. This gradient integral then
vanishes as well. Inner regularity gives absolute continuity
of the derivative measure and proves Theorem~\ref{thm:main}.

Lower-radius Lipschitz truncations appear in Haj{\l}asz and
Mal\'y \cite[Lemma~8]{HajlaszMaly2010}, and local/nonlocal
splitting with subtraction of an input mean appears in Saari
\cite[Section~3.1]{Saari2019}. The signed estimate above supplies
the gradient-only localization used here.

Section~\ref{sec:local} proves disk cancellation and the local estimates.
Section~\ref{sec:prices} gives scalar refinement and common prices; Section~\ref{sec:labels}
constructs the vector particles and spatial label law.
Section~\ref{sec:assembly} assembles the finite bands, and Section~\ref{sec:extension} passes to
Sobolev input and removes the singular part. An analytic appendix records the Sobolev
calculus, square localization, and vector-measure facts in the
forms used in the final step. \subsection{Conventions}


A constant denoted by $C$ may increase between occurrences.
All parameters, including cone apertures, finite direction nets and
geometric neighborhood sizes, are fixed absolute choices. In a local
lemma the explicitly listed cone and support constants are allowed
dependencies. No constant depends on the number of radius scales.

For a vector measure $\nu$, $|\nu|$ denotes its total variation.
For a real number $a$, $a_+=\max(a,0)$.
A supremum of nonnegative quantities over an empty set is zero.
When $\nabla u_t$ is written, the radius is held fixed.
We use the classical strong $L^p$ maximal inequality for $p>1$
\cite[(1.2)]{Kinnunen1997}; a general account is given in
\cite[Chapter~I]{Stein1970}. The analytic appendix records the
remaining standard inputs in the forms used here.
